\chapter{Conclusion}
\label{Chapter6}

\section{Research Summary}
\label{sec:conclusion-summary}

\subsection{Research Objectives}

The research objectives were achieved within the evaluated scope. The study identified
hybrid-cloud operational gaps, proposed and implemented a SysSecOps process, and developed
a common finding representation with an explainable risk-scoring mechanism. The
evaluation established internal correctness and practical feasibility, but not
large-scale reliability or general effectiveness across cloud platforms.

\subsection{RQ1: Decision Conformance and Enforcement}

RQ1 was answered positively under the declared conditions. All $106$ pre-registered runs
produced their expected decision, and all six enforcement invariants held whenever
applicable. However, degradation testing exposed a fail-open weakness in which missing
assessment evidence could move a review decision to pass; the decision policy therefore
requires an explicit assurance penalty for unavailable components.

\subsection{RQ2: Cross-Boundary Enforcement}

RQ2 was answered positively for workflow enforcement, with bounded detection coverage.
The on-premises path attained all seven operational checkpoints, including live apply,
verification, and idempotency. Shared detection increased from three to ten of eleven
weakness classes after remediation, showing that a common risk engine requires comparable
assessment coverage and finding vocabulary on both sides of the hybrid boundary.

\subsection{RQ3: Unified Risk Evaluation}

RQ3 was partially answered. On $324$ held-out Python samples, the learned component
achieved an accuracy of $0.9259$, an F1-score of $0.9240$, and a ROC--AUC of $0.9233$.
Its probability contributed successfully to the unified score, but the result remains
limited to detector-visible Python vulnerabilities and does not establish generalisation
to other languages or assessment environments.

\section{Research Contributions}
\label{sec:conclusion-contributions}

This research makes three principal contributions.

\begin{enumerate}
    \item \textbf{A unified SysSecOps operational process.} This contribution was
    achieved by connecting infrastructure validation, security assessment, approval,
    deployment, and monitoring through one auditable workflow for cloud and on-premises
    targets. The $106$-run campaign and live-target arm showed that the implemented
    process reaches and enforces its declared decisions.

    \item \textbf{A common risk and finding mechanism.} This contribution was achieved
    by normalising heterogeneous findings into a shared vocabulary, deduplicating them,
    and combining severity, operational context, and learned probability in an
    explainable score. Paired fixtures exposed cross-boundary coverage gaps and guided
    remediation from three to ten commonly detected weakness classes.

    \item \textbf{A reproducible prototype and evaluation method.} This contribution was
    achieved through version-controlled scenarios, pre-declared expectations, persisted
    machine-readable evidence, enforcement invariants, live-target verification, and a
    labelled code-risk evaluation. These artifacts allow decisions and experimental
    conclusions to be reproduced and audited.
\end{enumerate}

\section{Future Work}
\label{sec:conclusion-future-work}

Future work follows directly from the limitations identified by the evaluation.

\begin{enumerate}
    \item \textbf{Close cross-boundary coverage gaps.} Because one cloud-side secret
    weakness remains unnamed and the private path lacks an equivalent learned signal,
    future work should add source-appropriate detection while preserving the common
    finding.

    \item \textbf{Strengthen external validity and scale.} Because the fixtures were
    authored within the study and the live-target arm used one host, evaluation should use
    independent fixtures, heterogeneous prior states, concurrent targets, induced
    failures, larger repositories, and incremental scanning.

    \item \textbf{Broaden and validate the learned model.} Because the model is limited
    to detector-visible Python samples, it should be evaluated with larger multi-project
    datasets, project-aware splits, repeated cross-validation, external holdout sets, and
    language-specific models.

    \item \textbf{Calibrate operational value.} Because the score weights and thresholds
    are not organisation-specific and no manual baseline was measured, future studies
    should use expert and incident data to calibrate the score and compare the process
    with manual operations and ungated infrastructure automation.
\end{enumerate}

\section{Concluding Remarks}
\label{sec:conclusion-remarks}

This research presents a structured approach to coordinating infrastructure management,
software delivery, security assessment, and operational monitoring in a hybrid cloud
environment. Its central proposition is that security should function as an explicit,
auditable operational decision rather than as an isolated activity after deployment.
The proposed process, unified scoring mechanism, and prototype provide a foundation for
this integration. The pre-registered campaign showed that the implemented gate decides
and enforces as specified; it also showed that the specification itself permits a change
to be released when evidence is missing, and that a shared decision function does not by
itself produce equal treatment across a hybrid boundary when detection coverage differs.
Those two findings are as much a result of this work as the conformance figures, and
they define the agenda for further validation. With broader evaluation, calibration, and
platform coverage, the approach can provide a basis for more consistent and traceable
SysSecOps practice across heterogeneous cloud infrastructure.

\endinput